\chapter{Stage 8 Robustness at Pre-Approved Strengths}
\label{chap:s2-4}

\section{Goal}

Does a verdict, and the region that explains it, survive the kind of image degradation a site camera
produces? Stage~8 had not been built. This phase builds it, with the perturbation strengths fixed and
approved before any run, so they could not be tuned against results.

\section{What Was Done}

Five perturbations, applied one at a time, never combined:

\begin{center}
\begin{tabular}{lp{9.2cm}}
\toprule
\textbf{perturbation} & \textbf{strength (approved 2026-11-10)} \\
\midrule
blur       & Gaussian, $\sigma = 1\%$ of the image's long side \\
low light  & $-2$ EV, applied in linear light, not gamma-encoded values \\
occlusion  & a mean-luminance grey block over the centre 50\% $\times$ 50\% of the image \\
contrast   & factor 0.4 \\
desaturate & ITU-R 601-2 luma \\
\bottomrule
\end{tabular}
\end{center}

The subset is 1,600 Stage~2 rows (800 YES, 800 NO, at least 40 per rule), seed 42; each row is asked again
under each perturbation with greedy decoding. Per perturbation the stage records the new verdict, whether it
matches the baseline, the region, and the baseline-to-perturbed centroid distance over the image diagonal.
A surviving verdict whose region moved more than 10\% of the diagonal is flagged --- the same cut-off Stage~6
now uses. For occlusion it also records how much of the annotated hazard the block hides, so the case where
the block covers the hazard itself is reported separately. The strengths are part of the run fingerprint.

\section{Results}

\begin{center}
\begin{tabular}{lrr}
\toprule
\textbf{perturbation} & \textbf{verdict survives} & \textbf{survived, region moved $>10\%$} \\
\midrule
blur       & 70.4\% [68.2, 72.6] & 47.7\% ($n=459$) \\
low light  & 90.7\% [89.2, 92.1] & 13.8\% ($n=884$) \\
occlusion  & 80.2\% [78.2, 82.1] & 35.1\% ($n=643$) \\
contrast   & 90.4\% [88.9, 91.8] & 14.8\% ($n=865$) \\
desaturate & 86.6\% [84.9, 88.2] & 20.9\% ($n=789$) \\
\bottomrule
\end{tabular}
\captionof{table}{Stage 8, A40, 4-bit, $n=1{,}600$ per perturbation. The drift column is over rows that
survived and named a region both times.}
\label{tab:s24}
\end{center}

Occlusion that hides at least 10\% of the annotated hazard leaves the verdict standing 67.9\% of the time
($n=290$); occlusion that hides less leaves it standing 85.1\% ($n=141$). That is the direction expected if
the model's verdict depends on the hazard rather than on the scene around it.

\section{Standard vs.\ Non-Standard Choices}

\textbf{Standard.} One-factor-at-a-time perturbation; survival rate per perturbation, never pooled.

\textbf{Non-standard.} Strengths fixed and approved before any run; low light simulated in linear light;
the occluder's coverage of the annotated hazard recorded per row, so occlusion results separate ``hid the
evidence'' from ``hid something else''.

\section{Implications}

Mild blur is the most damaging change: three verdicts in ten flip, and when the verdict survives, the
region moves in almost half the cases. A surviving verdict is not a surviving explanation --- the paper
should report both columns.
